\documentclass[12pt]{article}
\title{Absolute acceleration in Galilean relativity}
\author{Hans Halvorson}
\date{\today}
\begin{document}

I take acceleration in Galilean relativity to be a paradigm case of an
absolute quantity. If we want to know which quantities are absolute in
other theories, e.g.\ General Relativity, then we had better get clear
on what exactly it means to say that acceleration is absolute in
Galilean spacetime.

There are (at least) two equivalent ways of describing Galilean
spacetime: as an affine space with some metric structure, and as a
manifold with a connection and some metric structure.

Recall that an \emph{affine space} is an object $a:A\times V\to A$,
where $A$ is a set, $V$ is a vector space, and $a$ is a




\end{document}


%%% Local Variables:
%%% mode: latex
%%% TeX-master: t
%%% End:
